% =====================================================================
%  RESUMEN (castellano) · una página · 200-300 palabras
%  Debe decir, por este orden: problema y contexto (2-3 frases), qué se ha
%  hecho y cómo, qué se ha obtenido (con cifras) y qué aporta.
%  No es una introducción: el lector tiene que saber el RESULTADO.
%  Resum y Abstract dicen lo mismo en su idioma, cada uno en su página.
% =====================================================================
\begin{otherlanguage}{spanish}
\chapter*{\resumenname}
\pdfbookmark[0]{\resumenname}{resumen}

% Ejemplo de estructura (sustitúyelo por tu propio resumen):
Reservar un espacio de estudio en el campus obliga hoy a consultar varias
fuentes y no permite saber si una sala está libre en ese momento. Este trabajo
presenta \emph{SalaLibre}, una aplicación móvil y un servicio web que muestran
la ocupación de las salas en tiempo real y permiten reservarlas en dos pasos.

La solución se ha desarrollado con una metodología iterativa en tres
incrementos. La aplicación Android se comunica con una API REST que guarda las
reservas en una base de datos relacional y notifica los cambios a los clientes.
Se han definido doce requisitos funcionales y cinco no funcionales, y cada uno
se ha validado con, al menos, un caso de prueba.

Las pruebas de carga muestran que, con caché, la latencia media se mantiene por
debajo de \qty{100}{ms} hasta 400 usuarios concurrentes, mientras que la versión
inicial superaba ese umbral a partir de 200. Una prueba de usabilidad con ocho
estudiantes obtuvo una puntuación SUS media de 81 sobre 100. Los datos de este ejemplo son ficticios.

\vspace{1em}
\noindent\textbf{\clavesname:} \tfgclavesES
\end{otherlanguage}
\cleardoublepage
